% ---------------------------------------------------------------
% seccoes/aquisicao.tex — Aquisição de conhecimento
% ---------------------------------------------------------------

\chapter{Aquisição de conhecimento}\label{sec:3}

A base de conhecimento foi construída de forma iterativa, a partir de quatro tipos de contributo: o perito, a regulamentação rodoviária, a literatura técnica e científica e o conhecimento da equipa (Figura~\ref{fig:2}). Cada fonte de conhecimento confirmou, corrigiu ou acrescentou critérios de decisão, e cada regra nova foi validada em cenários e testes antes de ser aceite.

\begin{figure}[H]
\centering
\begin{tikzpicture}[
    node distance=6mm and 9mm,
    caixa/.style={draw=corEscura, line width=0.8pt, rounded corners=3pt, fill=corFundo,
                  align=center, font=\small, minimum height=9mm, text width=50mm},
    fonte/.style={caixa, text width=34mm},
    seta/.style={-{Stealth[length=2.4mm]}, line width=0.8pt, draw=corEscura}
]
\node[caixa] (aq) {Aquisição de conhecimento\\do domínio};
\node[caixa, below=of aq] (mod) {Modelo conceptual: factos,\\eventos e prioridades};
\node[caixa, below=of mod] (form) {Formalização das regras};
\node[caixa, below=of form] (val) {Validação de cenários e testes};
\node[caixa, right=10mm of val, text width=34mm] (rev) {Revisão da base\\de conhecimento};
\node[fonte, left=12mm of aq, yshift=9mm] (f1) {Literatura académica\\e técnica};
\node[fonte, left=12mm of aq, yshift=-8mm] (f2) {Perito do domínio\\Prof.\ Gil Lopes};
\draw[seta] (f1.east) -- (aq.west);
\draw[seta] (f2.east) -- (aq.west);
\draw[seta] (aq) -- (mod);
\draw[seta] (mod) -- (form);
\draw[seta] (form) -- (val);
\draw[seta] (val) -- (rev);
\draw[seta] (rev.north) |- node[pos=0.25, right, font=\scriptsize, text=corCinza]{iteração} (aq.east);
\end{tikzpicture}
\caption{Processo de aquisição, validação e atualização da base de conhecimento}\label{fig:2}
\end{figure}

\section{Fontes de conhecimento}\label{sec:3.1}

A primeira fonte considerada corresponde à pesquisa bibliográfica e técnica, utilizada para aprofundar conceitos relacionados com os sistemas periciais, sensores, segurança rodoviária e a avaliação de situações de risco. Esta pesquisa permitiu também compreender melhor determinados indicadores relevantes para o sistema, como o TTC, utilizado para relacionar a distância a um objeto com a velocidade de aproximação.

A segunda fonte diz respeito ao conhecimento desenvolvido pelo grupo, nomeadamente à identificação dos elementos considerados importantes para a tomada de decisão, às possíveis relações entre estes elementos e as ações que o sistema poderia propor. Este conhecimento foi utilizado para construir uma primeira versão das regras a aplicar e, simultaneamente, permitiu identificar dúvidas a esclarecer com o perito.

Por último, a fonte de conhecimento mais importante foi o contacto com o perito, não só para adquirir novas informações e conhecimento, mas também para rever e questionar algumas decisões, inicialmente, definidas pelo grupo. Esta fonte de conhecimento permitiu esclarecer aspetos do domínio relacionados, por exemplo, com critérios de criticidade e prioridades. Desta forma, a base de conhecimento é o resultado da combinação entre a pesquisa, o conhecimento desenvolvido pelo grupo e o conhecimento especializado, como apresentado, resumidamente, pela Tabela~\ref{tab:2}.

\begin{tabela}{L{\dimexpr0.2026\linewidth-2\tabcolsep\relax}L{\dimexpr0.3588\linewidth-2\tabcolsep\relax}L{\dimexpr0.4367\linewidth-2\tabcolsep\relax}}
\caption{Fontes de conhecimento e respetivo contributo}\label{tab:2}\\
\linhacab \cab{\textbf{Tipo de fonte}} & \cab{\textbf{Fonte}} & \cab{\textbf{Contributo para a base de conhecimento}} \\
\endfirsthead
\linhacab \cab{\textbf{Tipo de fonte}} & \cab{\textbf{Fonte}} & \cab{\textbf{Contributo para a base de conhecimento}} \\
\endhead
Perito & Prof. Doutor Agostinho Gil Teixeira Lopes, ISEP (Secção~\ref{sec:3.2}) & Criticidade a partir da velocidade, da distância e da capacidade de travagem; número de leituras antes de agir; níveis de risco e desempate; validade dos factos; comportamento em caso de falha \\
\rowcolor{corFundo}Legislação & Código da Estrada [7] & Paragem no vermelho e no STOP, cedência a peões, circulação pela direita, ultrapassagem pela esquerda, velocidade adequada ao espaço livre e visível \\
Guia técnico & Prevenção Rodoviária Portuguesa [8] & Distância de segurança e condições climatéricas \\
\rowcolor{corFundo}Literatura científica & Vogel [9] & TTC como indicador de risco de colisão \\
Literatura científica & Yeong et al. [10]; Wei et al. [11]; Castanedo [12] & Escolha dos sensores e fusão de dados \\
\rowcolor{corFundo}Material da unidade curricular & Materiais de ENGCIA [13] & \emph{Drools}, raciocínio temporal, lógica difusa \\
Equipa & Análise do domínio e dos cenários & Primeira versão das regras e questões a colocar ao perito \\
\end{tabela}

\section{O perito}\label{sec:3.2}

O perito selecionado para o desenvolvimento do presente projeto é o Professor Doutor Agostinho Gil Teixeira Lopes, atualmente Professor Adjunto do Departamento de Engenharia Informática do Instituto Superior de Engenharia do Porto (ISEP), desde outubro de 2024 [1, 2]. A escolha deste perito é fundamentada pela formação académica, experiência profissional e atividade de investigação nas áreas diretamente relacionadas com o domínio do trabalho, nomeadamente a robótica autónoma, a perceção sensorial, a automação e os sistemas de controlo.

No que respeita à formação académica, o perito concluiu, em 2002, a licenciatura em Informática e Matemáticas Aplicadas na Universidade Portucalense Infante D. Henrique e, em 2006, o doutoramento em Engenharia Informática na \emph{Cranfield University}, no Reino Unido, com especialização em \emph{Automation, Control and Robotics}. A investigação desenvolvida no âmbito da tese apresentada, intitulada \emph{Arc-Based Sensing in Narrow Groove Pipe Welding}, teve como foco a utilização de informação sensorial em processos de automação e controlo. Seguiram-se, posteriormente, dois projetos de pós-doutoramento em \emph{Cranfield} (2006–2008) [3].

Relativamente ao percurso profissional, conta com a atividade de investigação na Universidade do Minho, entre 2008 e 2013, bem como com funções de Professor Auxiliar na Universidade da Maia, entre 2007 e 2024, e de Professor Auxiliar Convidado na Universidade do Minho, entre 2013 e 2024 [1]. Ainda no decorrer da carreira profissional, participou em projetos de investigação, publicou 76 artigos científicos, 42 dos quais indexados na \emph{Scopus} [1], e orientou projetos académicos em diferentes áreas da robótica e da inteligência artificial. Atualmente, a área de investigação para a qual contribui centra-se na robótica móvel autónoma inteligente, com temas como visão computacional, sistemas multissensoriais, localização, navegação e controlo de plataformas autónomas.

Neste sentido, a escolha deste perito justifica-se, sobretudo, pela proximidade entre as respetivas áreas de especialização e os desafios associados ao sistema pericial proposto. Embora o projeto não tenha como finalidade o desenvolvimento de um sistema de condução autónoma para utilização num veículo real, a representação do conhecimento necessária à análise dos cenários simulados depende da compreensão de como a informação deve ser interpretada e utilizada na avaliação de situações e na definição das ações realizadas.

O contributo do perito foi particularmente pertinente na discussão de aspetos como, por exemplo, a relação entre velocidade, distância e capacidade de travagem, a identificação de situações críticas, o número de leituras necessário para confirmar uma deteção e o tratamento de informação contraditória. Estas questões encontram-se diretamente relacionadas com a aquisição e representação do conhecimento, uma vez que permitem fundamentar as condições, exceções e prioridades consideradas na construção da base de conhecimento. Assim, a participação do perito permitiu complementar a pesquisa realizada pela equipa, questionar pressupostos, inicialmente, definidos e orientar a formalização de conhecimento relevante para um sistema pericial cujas decisões possam ser compreendidas e analisadas pelos utilizadores.

No entanto, uma vez que o perito não é especialista em legislação rodoviária, nem em engenharia e homologação automóvel, o conhecimento normativo foi consultado diretamente no Código da Estrada.

\section{Sessões com o perito}\label{sec:3.3}

O contacto com o perito decorreu em três momentos (Tabela~\ref{tab:3}). As duas primeiras sessões tiveram como objetivo adquirir conhecimento e esclarecer dúvidas relacionadas com a construção da base de conhecimentos, nomeadamente sobre as variáveis avaliar, as condições que pudessem alterar uma decisão, as exceções e o tratamento de conflitos. A terceira sessão teve como objetivo a apresentação da solução implementada e à respetiva apreciação do perito. O resumo das perguntas e respostas, da sessão 1 e sessão 2, encontra-se no Anexo A.

\begin{tabela}{L{\dimexpr0.1759\linewidth-2\tabcolsep\relax}L{\dimexpr0.2740\linewidth-2\tabcolsep\relax}L{\dimexpr0.2740\linewidth-2\tabcolsep\relax}L{\dimexpr0.2740\linewidth-2\tabcolsep\relax}}
\caption{Sessões com o perito}\label{tab:3}\\
\linhacab \cab{} & \cab{\textbf{Sessão 1 (S1)}} & \cab{\textbf{Sessão 2 (S2)}} & \cab{\textbf{Sessão 3 (S3)}} \\
\endfirsthead
\linhacab \cab{} & \cab{\textbf{Sessão 1 (S1)}} & \cab{\textbf{Sessão 2 (S2)}} & \cab{\textbf{Sessão 3 (S3)}} \\
\endhead
\textbf{Data e formato} & 1 de outubro de 2026 Reunião via Microsoft Teams & 2 de outubro de 2026\par Email com questões & 24 de outubro de 2026\par Reunião via Microsoft Teams \\
\rowcolor{corFundo}\textbf{Objetivo} & Apresentar o projeto e discutir a respetiva implementação & Esclarecer as questões que ficaram em aberto & Apresentar o sistema em funcionamento e apreciação final do perito \\
\textbf{Temas} & Estrutura do sistema, representação do ambiente, factos, cenários e influência dos factos na decisão & Falhas, relevância dos sensores, criticidade e TTC, número de leituras, prioridades e desempate, memória de trabalho, dados contraditórios, cenários, explicações & Simulação e decisões em tempo real, construção de cenários na matriz de vias, funcionamento do motor de inferência e explicações \\
\rowcolor{corFundo}\textbf{Principais resultados} & Representar o percurso como uma matriz de vias e posições; centrar o risco na relação entre velocidade e distância; Teoria dos Jogos como evolução possível & Criticidade calculada pela distância de travagem; agir à primeira leitura em situação crítica e confirmar em situação normal; quatro níveis de risco; desempate determinístico & Apreciação muito positiva do sistema e, em particular, da matriz de vias \\
\end{tabela}

As indicações das duas primeiras sessões com o perito, estão refletidas no sistema através da implementação da matriz de vias, da interface na qual é possível inserir novos elementos na simulação, das regras aplicáveis aos peões, aos obstáculos e aos semáforos, assim como através do cálculo da distância de travagem, em detrimento de valores fixos para as distâncias e da implementação das regras de TTC que distinguem a reação imediata da confirmação por leituras sucessivas.

Por sua vez, na terceira sessão, a equipa apresentou a solução implementada e demonstrou o seu funcionamento através de uma simulação em tempo real. A reação do perito foi muito positiva. Mostrou-se bastante satisfeito com o funcionamento do sistema e com a forma como cada decisão é justificada, e valorizou de forma particular a matriz de vias, que permite criar e alterar cenários de forma rápida e interativa

A Tabela~\ref{tab:4} apresenta a forma como as principais indicações do perito foram representadas na base de conhecimento e o respetivo estado de implementação.

\begin{tabela}{L{\dimexpr0.3101\linewidth-2\tabcolsep\relax}L{\dimexpr0.5154\linewidth-2\tabcolsep\relax}L{\dimexpr0.1724\linewidth-2\tabcolsep\relax}}
\caption{Das indicações do perito à base de conhecimento}\label{tab:4}\\
\linhacab \cab{\textbf{Indicação do perito}} & \cab{\textbf{Representação no sistema}} & \cab{\textbf{Estado}} \\
\endfirsthead
\linhacab \cab{\textbf{Indicação do perito}} & \cab{\textbf{Representação no sistema}} & \cab{\textbf{Estado}} \\
\endhead
Representar a estrada como uma matriz & Estrada com 3 vias × blocos de 5 m, à qual é possível adicionar elementos no decorrer da simulação & Implementada \\
\rowcolor{corFundo}Criticidade por uma equação de velocidade, distância e capacidade de travagem; nunca travagem instantânea & Distância de paragem com tempo de reação de 0,5 s e desaceleração de 2,8 m/s² (conforto) ou 8,5 m/s² (emergência), usada em semáforos, STOP, peões, obstáculos, barreiras e buracos & Implementada \\
Situação crítica: agir à primeira leitura; normal: confirmar & TTC < 1,3 s numa leitura dispara a emergência;\par TTC < 2,5 s exige 3 leituras em 2 s & Implementada \\
\rowcolor{corFundo}Níveis de risco & Quatro níveis de risco e a decisão usa sete categorias de prioridade (Secção~\ref{sec:4.4}) & Implementada \\
Desempate: escolher a primeira de uma lista ordenada & Em empate, vence a proposta de menor velocidade e, depois, a primeira registada & Implementada \\
\rowcolor{corFundo}Esquecer o que sai do campo de deteção; um sinal é uma mudança de estado & Perceções expiram no caso de não haver novas leituras (250 ms a 1,5 s) & Implementada \\
Acrescentar a passadeira com peão e elementos fora de contexto & Cenário S7;\par é possível também inserir elementos com a simulação a correr & Implementada \\
\rowcolor{corFundo}As explicações devem estar alinhadas com as regras & O texto de cada explicação é a razão e a evidência registada pela regra que disparou & Implementada \\
Teoria dos Jogos & Não seguida; ver desenvolvimento futuro (Secção~\ref{sec:7.3}) & Futura \\
\end{tabela}

As indicações recolhidas junto do perito, em conjunto com as restantes fontes de conhecimento, foram formalizadas na base de conhecimento, como se descreve no capítulo seguinte. A apresentação final permitiu demonstrar a aplicação prática destas indicações e recolher uma apreciação qualitativa da solução desenvolvida.
